%!TEX TS-program = pdflatex
%!TEX TS-options = -shell-escape
% % % % %   Die folgenden Zeilen müssen ihre Zeilennummern 4 und 5 behalten !!!    % % % % %
\newcommand{\printpraesenzlsg}{false}
\newcommand{\printloesungen}{false}
\newcommand{\printbewertungen}{false}
% % % % %   \newcommand{\printloesungen}{false}                                    % % % % %
\newcommand{\blattnummer}{1}
%\newcommand{\abgabetermin}{\textcolor{red}{bis 11.04.2022, 09:00 Uhr}}
\documentclass[%
  paper=a4,
  fontsize=10pt,
  ngerman
  ]{scrartcl}

% Basics für Codierung und Sprache
% ===========================================================
  \usepackage{scrtime}
  \usepackage{etex}
  \usepackage{shellesc}
  \usepackage[final]{graphicx}
  \usepackage[utf8]{inputenc}
  \usepackage{babel}
  \usepackage[german=quotes]{csquotes}
  \usepackage{datetime2}
  \usepackage{ifthen}
  \usepackage{dingbat}
% ===========================================================

%% Fonts und Typographie
%% ===========================================================
  \usepackage{charter}
  \usepackage{nimbusmononarrow}
  \usepackage[babel=true,final,tracking=smallcaps]{microtype}
  \DisableLigatures{encoding = T1, family = tt* }
  \usepackage{ellipsis}
% ===========================================================

% Farben
% ===========================================================
  \usepackage[usenames,x11names,final,table]{xcolor}
  \definecolor{fbblau}{HTML}{3078AB}
  \definecolor{blackberry}{rgb}{0.53, 0.0, 0.25}
% ===========================================================

% Mathe-Pakete und -Einstellungen
% ===========================================================
  \usepackage{mathtools}
  \usepackage{amssymb}
  \usepackage[bigdelims]{newtxmath}
  \allowdisplaybreaks
  \usepackage{bm}
  \usepackage{wasysym}
% ===========================================================

% TikZ
% ===========================================================
  \usepackage{tikz}
  \usepackage{tikz-cd}
  \usepackage{pgf}
  \usetikzlibrary{matrix}
  \usetikzlibrary{positioning}
  \usetikzlibrary{arrows}
  \usetikzlibrary{shapes}
  \tikzset{>=Latex}
  \usepackage[mode=image]{standalone}
% ===========================================================

% Seitenlayout, Kopf-/Fußzeile
% ===========================================================
  \usepackage{scrlayer-scrpage}
  \usepackage[top=3cm, bottom=2.5cm, left=2.5cm, right=2cm]{geometry}  
  \pagestyle{scrheadings}
  \clearscrheadfoot 
  \setheadsepline{0.4pt}
  \setfootsepline{0.4pt}
  \setkomafont{pagehead}{\normalfont\footnotesize}
  \setkomafont{pagefoot}{\normalfont\footnotesize}
  \lohead[]{Klausur zur Vorlesung \textit{Mathematik I - Theoretische Grundlagen der Informatik} -- Wintersemester 2022/2023}
  \rohead[]{%
    \ifthenelse{\boolean{loesungen}}{\textbf{Lösungen zu Blatt \blattnummer}}{Blatt \blattnummer}}
  %\rofoot[Version: \DTMnow]{%
  %  \ifthenelse{\boolean{loesungen}}{Version: \DTMnow}{}}
  %\lofoot[\jobname.tex]{%
  %  \ifthenelse{\boolean{loesungen}}{\jobname.tex}{}}
  \cfoot[\thepage]{\thepage}
  \raggedbottom
  \usepackage{setspace}          
  \linespread{1.2}
  \setlength{\parindent}{0pt}
  \setlength{\parskip}{0.5\baselineskip}
  \usepackage[all]{nowidow}
  \usepackage{lscape}
% ===========================================================

% Hyperref
% ===========================================================
  \usepackage[%
    hidelinks,
    pdfpagelabels,
    bookmarksopen=true,
    bookmarksnumbered=true,
    linkcolor=black,
    urlcolor=fbblau,
    plainpages=false,
    pagebackref,
    citecolor=black,
    hypertexnames=true,
    pdfborderstyle={/S/U},
    linkbordercolor=fbblau,
    colorlinks=false,
    backref=false]{hyperref}
  \hypersetup{final}
  \makeatletter
  \g@addto@macro{\UrlBreaks}{\UrlOrds}
  \makeatother
% ===========================================================

% Listen und Tabellen
% ===========================================================
  \usepackage{multicol}
  \usepackage{multirow}
  \usepackage[shortlabels]{enumitem}
  \setlist[enumerate]{font=\bfseries,itemsep=0.25\baselineskip}
  \setlist[itemize]{label=$\triangleright$,itemsep=-0.25\baselineskip}
  \usepackage{tabularx}
  \usepackage{array}
  \usepackage{listings}
  \lstset{breaklines=true,
    basicstyle=\ttfamily,
    frame=tblr,
    language=c,
    numbers=left}
  \usepackage{include/mips}

\definecolor{lightgreen}{rgb}{0,0.7,0}
\definecolor{gray}{rgb}{0.5,0.5,0.5}
\definecolor{mauve}{rgb}{0.58,0,0.82}

\lstdefinestyle{mips}{ %
  language=[mips]Assembler,       % the language of the code
  keywordstyle=\color{blue},          % keyword style
  keywordstyle=[2]\color{magenta},     % keyword style
  keywordstyle=[3]\color{red},          % keyword style
  commentstyle=\color{lightgreen},       % comment style
  stringstyle=\color{mauve},         % string literal style
  escapeinside={\%*}{*)},            % if you want to add a comment within your code
  morekeywords={*,...}               % if you want to add more keywords to the set
}

\newcommand{\mipsinline}[1]{\lstinline[style=mips,columns=fixed]{#1}}
\newcommand{\mipsregister}[1]{\mipsinline{\$#1}}
  \newcommand{\qed}{\hfill \ensuremath{\Box}}
  \newcommand{\cmd}[1]{\texttt{#1}}
% ===========================================================

% Sonstiges
% ===========================================================
  \usepackage[textsize=small,color=Red1!80!OrangeRed1!80]{todonotes}
  \usepackage{lipsum}
% ===========================================================

% Shortcuts and Math Commands
% ===========================================================
  \newcommand{\BB}{\mathbb{B}}      % Bool
  \newcommand{\CC}{\mathbb{C}}      % komplexe Zahlen
  \newcommand{\NN}{\mathbb{N}}      % nat. Zahlen
  \newcommand{\QQ}{\mathbb{Q}}      % rat. Zahlen
  \newcommand{\RR}{\mathbb{R}}      % reelle Zahlen
  \newcommand{\ZZ}{\mathbb{Z}}      % ganze Zahlen
  \newcommand{\oh}{\mathcal{O}}      % Landau-O
  \newcommand{\ol}[1]{\overline{#1}}    % überstreichen
  \newcommand{\wt}[1]{\widetilde{#1}}    % überschlängeln
  \newcommand{\wh}[1]{\widehat{#1}}    % überdachen
  \newcommand{\entspricht}{\mathrel{\widehat{=}}}
  
  \DeclarePairedDelimiter{\absolut}{\lvert}{\rvert}    % Betrag
  \DeclarePairedDelimiter{\ceiling}{\lceil}{\rceil}    % aufrunden
  \DeclarePairedDelimiter{\Floor}{\lfloor}{\rfloor}    % aufrunden
  \DeclarePairedDelimiter{\Norm}{\lVert}{\rVert}      % Norm
  \DeclarePairedDelimiter{\sprod}{\langle}{\rangle}    % spitze Klammern
  \DeclarePairedDelimiter{\enbrace}{(}{)}          % runde Klammern
  \DeclarePairedDelimiter{\benbrace}{\lbrack}{\rbrack}  % eckige Klammern
  \DeclarePairedDelimiter{\penbrace}{\{}{\}}        % geschweifte Klammern
  \newcommand{\Underbrace}[2]{{\underbrace{#1}_{#2}}}   % bessere Unterklammerungen
  
  % Kurzschreibweisen für Faule und Code-Vervollständigung
  \newcommand{\abs}[1]{\absolut*{#1}}
  \newcommand{\ceil}[1]{\ceiling*{#1}}
  \newcommand{\flo}[1]{\Floor*{#1}}
  \newcommand{\no}[1]{\Norm*{#1}}
  \newcommand{\sk}[1]{\sprod*{#1}}
  \newcommand{\enb}[1]{\enbrace*{#1}}
  \newcommand{\penb}[1]{\penbrace*{#1}}
  \newcommand{\benb}[1]{\benbrace*{#1}}
  \newcommand{\stack}[2]{\makebox[1cm][c]{$\stackrel{#1}{#2}$}}

% Übungszettel-Krams
% ===========================================================
  \usepackage[tikz]{mdframed}
  \newmdenv[%
    backgroundcolor = gray!20,
    linewidth=0pt,
    skipabove = 0.5cm,
    skipbelow = 0cm
  ]{notes}
  
  % Lösungsausgabe
  \newboolean{praesenzlsgn}
  \setboolean{praesenzlsgn}{\printpraesenzlsg}
  \newboolean{loesungen}
  \setboolean{loesungen}{\printloesungen}
  % Bewertungsschema-Ausgabe
  \newboolean{bewertungen}
  \setboolean{bewertungen}{\printbewertungen}

  \newcommand{\iforiginal}[1]{\ifthenelse{\boolean{praesenzlsgn}}{}{#1}}
  \newcommand{\ifpraesenzlsg}[1]{\ifthenelse{\boolean{praesenzlsgn}}{#1}{}}
  \newcommand{\ifloesung}[1]{\ifthenelse{\boolean{loesungen}}{#1}{}}
  \newcommand{\ifbewertung}[1]{\ifthenelse{\boolean{bewertungen}}{#1}{}}

  \newcommand{\umbruchoriginal}{\iforiginal{\newpage}}        % Seitenumbruch nur im Original
  \newcommand{\umbruchpraesenzlsg}{\ifpraesenzlsg{\newpage}}  % Seitenumbruch nur in Lösung d. Präsenzaufgabe
  \newcommand{\umbruchlsg}{\ifloesung{\newpage}}              % Seitenumbruch nur in Lösung
  \newcommand{\umbruchbewertung}{\ifbewertung{\newpage}}      % Seitenumbruch nur in Lösung mit Bewertungsschema
  \newcommand{\umbruchnurpraesenz}{\ifthenelse{\boolean{loesungen}}{}{\umbruchpraesenzlsg}}
  \newcommand{\umbruchnurlsg}{\ifthenelse{\boolean{bewertungen}}{}{\umbruchlsg}}

  \usepackage{comment}
  \newenvironment{praesenzlsg}%
    {\ifthenelse{\boolean{praesenzlsgn}}{\textbf{--- Präsenzaufgabe Anfang ---} \mbox{} \newline}{}}%
	{\ifthenelse{\boolean{praesenzlsgn}}{\mbox{} \newline \textbf{--- Präsenzaufgabe Ende ---}}{}}
  \newenvironment{loesung}%
    {\ifthenelse{\boolean{loesungen}}{\textbf{--- Lösung Anfang ---} \mbox{} \newline}{}}%
    {\ifthenelse{\boolean{loesungen}}{\mbox{} \newline \textbf{--- Lösung Ende ---} \ifthenelse{\boolean{bewertungen}}{}{\umbruchlsg}}{}}
  \newenvironment{bewertung}%
    {\ifthenelse{\boolean{bewertungen}}{\textbf{--- Bewertungsschema Anfang ---} \mbox{} \newline}{}}%
    {\ifthenelse{\boolean{bewertungen}}{\mbox{} \newline \textbf{--- Bewertungsschema Ende ---} \umbruchlsg}{}}

  \ifthenelse{\boolean{praesenzlsgn}}{}{\excludecomment{praesenzlsg}}
  \ifthenelse{\boolean{loesungen}}{}{\excludecomment{loesung}}
  \ifthenelse{\boolean{bewertungen}}{}{\excludecomment{bewertung}}

  % Aufgabennummerierung
  \newcounter{aufgabennummer} 
  \setcounter{aufgabennummer}{1}
  \newcommand{\aufgabe}[1]{\textbf{Aufgabe \blattnummer.\theaufgabennummer} \stepcounter{aufgabennummer} \hfill (#1 Punkte)}
  %\newcommand{\aufgabetitel}[2]{\iforiginal{\vspace{0.5cm}} \textbf{Aufgabe \blattnummer.\theaufgabennummer \ (#2)} \stepcounter{aufgabennummer} \hfill (#1 Punkte)}
  \newcommand{\aufgabetitel}[2]{\iforiginal{\vspace{0.5cm}} \textbf{Aufgabe \theaufgabennummer \ (#2)} \stepcounter{aufgabennummer} \hfill (#1 Punkte)}
  \newcommand{\praesenztitel}[1]{\iforiginal{\vspace{0.5cm}} \textbf{Aufgabe \blattnummer.\theaufgabennummer \ (#1)} \stepcounter{aufgabennummer} \hfill \smallpencil}

\newcommand{\danger}{\raisebox{0.5mm}{\fontencoding{U}\fontfamily{futs}\selectfont\char 66\relax}}
\newcommand{\ddanger}[1]{{\danger} \textbf{#1} {\danger}}

% Rechnerstrukturen-Krams
% ===========================================================
  \newcommand{\Band}{\ensuremath{\mathbin{\texttt{AND}}}}
  \newcommand{\Bor}{\ensuremath{\mathbin{\texttt{OR}}}}
  \newcommand{\Bnot}{\ensuremath{\mathbin{\texttt{NOT}}}}
  \newcommand{\Bnand}{\ensuremath{\mathbin{\texttt{NAND}}}}
  \newcommand{\Bnor}{\ensuremath{\mathbin{\texttt{NOR}}}}
  \newcommand{\Bxor}{\ensuremath{\mathbin{\texttt{XOR}}}}
  \newcommand{\Bxnor}{\ensuremath{\mathbin{\texttt{XNOR}}}}
  \newcommand{\Bone}{\ensuremath{\mathbin{\texttt{ONE}}}}
  \newcommand{\Bzero}{\ensuremath{\mathbin{\texttt{ZERO}}}}
  \newcommand{\Binary}[1]{\ensuremath{\texttt{#1}}\xspace}
  \newcommand{\ieee}[3]{\fbox{\ensuremath{\texttt{#1}}}\fbox{\ensuremath{\texttt{#2}}}\fbox{\ensuremath{\texttt{#3}}}}
  \DeclareMathOperator{\DNF}{DNF}
  \DeclareMathOperator{\KNF}{KNF}
  \DeclareMathOperator{\xor}{\oplus}
  \newcommand{\dc}{\star} %don't-care-Wert
 

% Änderungen 2020: Hinweise auf Digitallehre angepasst; Rechnergeschichte entfernt; Blatt 1 und 2 zusammengefasst.
\begin{document}
\iforiginal{\thispagestyle{scrplain}
\vspace*{-3cm}
\begin{minipage}[t][1.1cm][c]{4.5cm}
  \includegraphics[width=4cm]{include/hwr-logo.png}
\end{minipage}
\hfill
\begin{minipage}[t][1.5cm][c]{8cm}
  \begin{center}
  \begin{footnotesize}
    \textsf{Prof. Dr.-Ing. Sebastian Schlesinger} \\[-0.1cm]
    \textsf{Fachbereich 2 - Duales Studium Wirtschaft \& Technik}
    %\url{https://www.hwr-berlin.de}
  \end{footnotesize}
  \end{center}
\end{minipage}
\hfill
%\begin{minipage}[t][1.6cm][c]{3.2cm}
%  \includegraphics[width=2.8cm]{include/ESlogo2}
%\end{minipage}

\vspace*{-0.3cm}
\begin{center} 
  \hrulefill \\[0.1cm]
  {\large Klausur} \\[0.15cm]
  {\huge \bfseries Mathematik I - Theoretische Informatik} \\[0.10cm]
  {HWR Berlin, Wintersemester 2024} \\[-0.4cm]
  \begin{tabular}{lcr}
    \hspace{0.3\textwidth}   & \hspace{0.3\textwidth} & \hspace{0.3\textwidth} \\
    Prof. Dr.-Ing. Sebastian Schlesinger   %&                        & "Ubung zur Klausurvorbereitung    \\ 
    %Blatt \blattnummer     % \abgabetermin          \\ 
  \end{tabular} \\[0.1cm]
  \hrulefill
\end{center}}

% \begin{notes} \small
% 	\textbf{Abgabetermine für Blatt 1:}
	
% 	Aufgaben 1.2/1.3: Montag, 11. April, 09:00 Uhr \\
% 	Aufgabe 1.4: Mittwoch, 20. April, 18:00 Uhr
% \end{notes}


\aufgabetitel{$5$}{Beweis}\\
Seien $A,B,C$ Mengen. Beweisen Sie \[(A\cap B)\cup C=A\cap (B\cup C)\Leftrightarrow C\subseteq A\]

\begin{loesung}
  \textbf{Beweis:}\\
\glqq $\Rightarrow$\grqq: Gelte $(A\cap B)\cup C=A\cap(B\cup C)$. Wir zeigen $C\subseteq A$.

Sei dazu $x\in C$. Dann ist $x\in (A\cap B)\cup C$ (Weil $x\in C$ kann ich was beliebiges dazu vereinigen) und nach Voraussetzung auch 
$x\in A\cap (B\cup C)$. Damit insbesonder $x\in A$, was zu zeigen war.

\glqq $\Leftarrow$\grqq:\\
Gelte $C\subseteq A$. Wir zeigen $(A\cap B)\cup C=A\cap(B\cup C)$.

Sei $x\in (A\cap B)\cup C$. Das ist "aquivalent zu $x\in (A\cap B)\vee x\in C$

$\Leftrightarrow (x\in A\wedge x\in B)\vee x\in C$

$\Leftrightarrow (x\in A\wedge x\in B)\vee (x\in A\wedge x\in C)$ (weil $C\subseteq A$ nach Voraussetzung ist $x\in C$ "aquivalent zu $x\in A\wedge x\in C$)

$\Leftrightarrow x\in (A\cap B)\cup x\in (A\cap C)$

$\leftrightarrow x\in A\cap(B\cup C)$ (Distributivgesetz), was zu zeigen war.

\qed

\end{loesung}

\aufgabetitel{$4$}{Symmetrische Differenz}\\
  Unter 
  \[A\triangle B:=(A\backslash B)\cup (B\backslash A)\] versteht man die \textit{symmetrische Differenz} der Mengen $A$ und $B$.
  \begin{enumerate}[(i)]
    \item Machen Sie sich anhand eines Venn-Diagramms klar, was unter der symmetrischen Differenz anschaulich zu verstehen ist.
    \item Beweisen Sie: $\forall A,B:A\triangle B=(A\cup B)\backslash(A\cap B)$.
  \end{enumerate}

\begin{loesung}
Venn-Diagramm zeichne ich jetzt mal nicht. Bleibt als "Ubung.

\textbf{Beweis:}\\
Sei $x\in A\triangle B=(A\backslash B)\cup (B\backslash A)$

$\Leftrightarrow (x\in A\wedge x\notin B)\vee (x\in B\wedge x\notin A)$

$\Leftrightarrow (x\in A\vee x\in B)\wedge (x\in A\vee x\notin A)\wedge (x\notin B\vee x\in B)\wedge (x\notin B\vee x\notin A)$ (Distributivit"at)

$\Leftrightarrow (x\in A\vee x\in B)\wedge (x\notin A\vee x\notin B)$ ($x\in A\vee x\notin A$ ist immer wahr, kann man also k"urzen)

$\Leftrightarrow x\in (A\cup B)\wedge \neq (x\in A\wedge x\in B)$ (De Morgan)

$\Leftrightarrow x\in(A\cup B)\backslash(A\cap B)$

\qed
\end{loesung}

\aufgabetitel{$4$}{Mengenbeweis}\\
Zeigen Sie für beliebige Mengen $A,B$:
\[A\cap (B\cup A)=A\]

\begin{loesung}
\glqq$\subseteq$\grqq:

Sei $x\in A\cap(B\cup A)$. Dann ist $x\in A$ (und irgendwas anderem), also $A\cap(B\cup A)\subseteq A$.

\glqq$\supseteq$\grqq:

Sei $x\in A$. Dann ist $x\in A\vee (x\in A\wedge x\in B)$.

$\Leftrightarrow (x\in A\vee x\in A)\wedge (x\in A\vee x\in B)$ (Distributivit"at)

$\Leftrightarrow x\in A\wedge (x\in A\vee x\in B)$

$\Leftrightarrow x\in A\cap(A\cup B)$

\qed
\end{loesung}
\aufgabetitel{$4$}{Potenzmengenbeweis}\\
Zeigen Sie für beliebige Mengen $A,B$:
\[\mathscr{P}(A\cap B)=\mathscr{P}(A)\cap\mathscr{P}(B)\]

\begin{loesung}
Sei $M\in\mathscr{P}(A\cap B)$

$\Leftrightarrow M\subseteq A\cap B$

$\Leftrightarrow \forall x\in M:x\in A\cap B$

$\Leftrightarrow \forall x\in M: x\in A\wedge x\in B$

$\Leftrightarrow M\subseteq A\wedge M\subseteq B$

$\Leftrightarrow M\in\mathscr{P}(A)\wedge M\in\mathscr{P}(B)$

$\Leftrightarrow M\in\mathscr{P}(A)\cap\mathscr{P}(B)$

\qed
\end{loesung}

\aufgabetitel{$11$}{Relationen}\\
Es sei $R=\{(1,1),(2,2),(3,3),(4,4),(1,2),(3,2),(2,4)\}$ auf der Menge $M=\{1,2,3,4\}$.
\begin{enumerate}[(i)]
\item Stellen Sie die Relation als Adjazenzmatrix dar.
\item Stellen Sie die Relation als Graph dar.
\item Ist die Relation reflexiv, irreflexiv, symmetrisch, antisymmetrisch oder transitiv?
\item Falls die Relation nicht transitiv ist, was m"usste man hinzuf"ugen, um sie mit m"oglichst wenigen zus"atzlichen Paaren / Pfeilen transitiv zu machen? Anders ausged"uckt - was ist die transitive H"ulle der Relation? (Allgemein ist die transitive H"ulle die kleinste $R$ enthaltende transitive Relation. \glqq Klein\grqq ist im Sinne der Inklusion (Teilmengenbeziehung) gemeint.)
\item Man erh"alt die reflexiv-transitive H"ulle auch rechnerisch wenn man $R_{trans}=\bigcup_{n\in\mathbb{N}}R^n$ bildet mit $R^0=Id$, der Identit"at, die jedes Element von $M$ mit sich selbst in Relation setzt und $R^{n+1}=R^n\circ R$. Man berechnet also diese $R^n$ und jeder Schritt f"ugt der Vereinigung der bereits von $R^0$ bis $R^{n-1}$ berechneten Relationen ggf. weitere Paare / Pfeile hinzu bis die Vereinigung schlie\ss lich transitiv wird. In der Praxis, bei endlichen Mengen, h"ort man auf, weitere Relationen hinzuzuf"ugen sobald der Vorgang station"ar wird, also nichts mehr hinzugef"ugt wird. Berechnen Sie die reflexiv-transitive H"ulle von $R$ auf dem beschriebenen Weg.
\item Ist $R_{trans}$ im vorliegenden Fall eine Ordnung?
\item Stellen Sie das Hasse-Diagramm von $R_{trans}$ dar. 
\end{enumerate}

\begin{loesung}
Die Matrix und den Graphen lass ich weg. Ist wirklich einfach.
Die Relation ist reflexiv und antisymmetrisch.
Um sie transitiv zu machen muss man die Paare $(1,4),(3,4)$ hinzuf"ugen.
Es ist $R^0=\{(1,1),(2,2),(3,3),(4,4)\}$, $R^1=R=R^0\cup R^1$ ($R$ ist schon reflexiv). Und $R^2=\{(1,2),(2,4),(3,2),(3,4),(1,4)\}$ und damit $R^0\cup R^1\cup R^2=R\cup\{(1,4),(3,4)\}=R_{trans}$ (die anderen Paare waren schon in $R$ und bei $R^3$ wird nichts mehr hinzugef"ugt.)
$R_{trans}$ ist eine Ordnung. Das Hasse-Diagramm lass ich weg. Im Diagramm ist unten die $1$ und die $3$, die eine Ebene weiter oben auf $2$ zeigen, die wiederum auf $4$ ganz oben zeigt.
\end{loesung}
\end{document}
